The Tinker-\/\-H\-P suite is made of 5 main programs that have been designed to run in a High Performance Computing (H\-P\-C) context, on G\-P\-Us and C\-P\-Us, leveraging M\-P\-I to run with multiple processes. All of these run based on various {\bfseries input files} and a {\bfseries command line}. Some input files are common to all the binaries\-:
\begin{DoxyEnumerate}
\item an {\bfseries xyz} file in the Tinker format representing a {\bfseries geometry} (through cartesian coordinates), {\bfseries atom types} of the atoms of the systems as well as their {\bfseries connectivity}.
\end{DoxyEnumerate}

{\itshape Below is an example of xyz file for a two water molecule system, oxygens are of atom type 101, hydrogens of type 88.}

``` 1 O 8.\-679662 7.\-087692 -\/0.\-696862 101 2 3 2 H 7.\-809455 6.\-755792 -\/0.\-382259 88 1 3 H 8.\-722232 6.\-814243 -\/1.\-617561 88 1 4 O -\/0.\-117313 8.\-244447 6.\-837616 101 5 6 5 H 0.\-216892 7.\-895445 6.\-050027 88 4 6 H 0.\-444268 7.\-826013 7.\-530196 88 4 ```


\begin{DoxyEnumerate}
\item a {\bfseries keyfile} containing various keywords describing the simulation such as the size of the box, the cutoffs on non-\/bonded simulations, the integrator for Molecular Dynamics, and also the {\bfseries parameter file}
\end{DoxyEnumerate}

{\itshape Below is an example of keyfile for a system using the charmm22 parameter file, with a cubic box size of edge 18.\-643 Angstroms and the verlet integrator.}

``` parameters charmm22 verbose a-\/axis 18.\-643 integrator verlet ```


\begin{DoxyEnumerate}
\item a {\bfseries parameter file} listing the parameters, for each atom type, of the force field used for the simulation.
\end{DoxyEnumerate}

For all programs of the Tinker-\/\-H\-P suite, the geometry needs to be included in the command line. To further indicate the programs that the simulation needs to take into account a specific keyfile, the keyfile needs to have the same prefix as the geometry or can be explicitely given in the command line with the prefix \char`\"{}-\/k keyfile\char`\"{}.

{\itshape Example of a command line for the analyze binary involving the watersmall.\-xyz geometry and the watersmall.\-key keyfile, on a single M\-P\-I process}

``` mpirun -\/np 1 analyze watersmall.\-xyz -\/k watersmall.\-key e ```

The 5 programs are listed below with a description of the command line necessary to run them. In the case of an incomplete command line or a missing input file an error message describing the missing point will be printed.


\begin{DoxyEnumerate}
\item {\bfseries analyze}\-: run single point energy given a trajectory file or a single frame.
\begin{DoxyItemize}
\item {\itshape command line}\-: E as in \char`\"{}\-Energy decomposition analysis\char`\"{} (see example above)
\end{DoxyItemize}
\item {\bfseries dynamic}\-: runs a molecular dynamics trajectory
\begin{DoxyItemize}
\item command line\-: a list of arguments has to be passed directly through the command line, respectively\-:
\begin{DoxyItemize}
\item number of steps of the desired trajectory
\item duration of a timestep (in femtoseconds)
\item frequency of output of frames (in picoseconds)
\item integer corresponding to the statistical ensemble to be sampled\-:
\begin{DoxyItemize}
\item 1\-: N\-V\-E ensemble
\item 2\-: N\-V\-T ensemble
\item 4\-: N\-P\-T ensemble
\end{DoxyItemize}
\item temperature (in Kelvin) for N\-V\-T and N\-P\-T ensembles
\item pressure (in Atmosphere) for N\-P\-T ensemble
\end{DoxyItemize}
\end{DoxyItemize}

{\itshape Example of a command line for the dynamic binary corresponding to 10000 time steps of 1fs in the N\-P\-T ensemble at 300\-K and 1 Atm, with frames printed out every picoseconds\-:}
\end{DoxyEnumerate}

``` mpirun -\/np 1 dynamic watersmall.\-xyz 10000 1 1 4 300 1 ```

More information about the features available with the {\bfseries dynamic} program, such as integrators, thermostats and barostats, can be found here.


\begin{DoxyEnumerate}
\item {\bfseries minimize}\-: runs a minimization of the potential energy of a system, starting from a geometry.
\begin{DoxyItemize}
\item {\itshape command line}\-: convergence criteria of the minimization algorithm
\end{DoxyItemize}
\end{DoxyEnumerate}

{\itshape Example of a command line for the minimize binary corresponding to a convergence threshold of 1 for R\-M\-S gradient\-:}

``` mpirun -\/np 1 minimize watersmall.\-xyz 1 ```

More information about the features available with the {\bfseries minimize} program, such as minimization algorithms and minimized geometry output can be found here.


\begin{DoxyEnumerate}
\item {\bfseries P\-I\-M\-D}\-: runs a Path Integral Molecular Dynamics trajectory to include nuclear quantum effects in the statistics.
\begin{DoxyItemize}
\item {\itshape command line}\-: the exact same inputs as the one of the {\bfseries dynamic} programs are necessary in the command line as described above. Several features specific to Path Integral such as the number of beads, as well as other implemented techniques to include nuclear quantum effects (the adaptive Quantum Thermal Bath) can be found in Quantum-\/\-H\-P.\-md.
\end{DoxyItemize}
\item {\bfseries bar}\-: compute free energy differences through the Bennett Acceptance Ratio estimator. In Tinker-\/\-H\-P, free energy differences are obtained in a post-\/processing step, in two stages. Given two keyfiles and two trajectories associated to two hamiltonians, first a $\ast$.bar file is produced with, for each trajectory, the potential energies associated to the two hamiltonians. Then, given the .bar file, one can run the self consistent procedure to get the final free energy difference through the bar estimator.
\begin{DoxyItemize}
\item {\itshape command line}\-:
\begin{DoxyItemize}
\item 1 to generate the .bar file, 2 to post-\/process it to get the bar estimation
\item for generation of the bar file, one has to indicate the two geometries with their keyfiles and the two temperatures
\end{DoxyItemize}

{\itshape Example of a command line for the bar binary to generate the .bar file corresponding to systema at temperature tempa and systemb at temperature tempb}
\end{DoxyItemize}
\end{DoxyEnumerate}

``` mpirun -\/np 1 bar 1 systema.\-xyz -\/k systema.\-key 300 systemb.\-xyz -\/k systemb.\-key 300 ```


\begin{DoxyItemize}
\item for estimation of the free energy difference, one has to indicate the bar file

{\itshape Example of a command line for the bar binary to get the free energy estimation based on the system.\-bar file}
\end{DoxyItemize}

``` mpirun -\/np 1 bar 2 system.\-bar ```


\begin{DoxyEnumerate}
\item {\bfseries testgrad}\-: compute gradients associated with a geometry. One can compute analytical gradients of the potential energy terms or numerical ones, on all the cartesian degrees of freedom.
\begin{DoxyItemize}
\item {\itshape command line}\-:
\begin{DoxyItemize}
\item Y to compute analytical gradients, N if it is not the case
\item Y to compute numerical gradients (through finite differences), N if not
\item when numerical gradients are requested, size of the step displacement in Angstroms
\item Y if breakdown by component is requested, N if not
\end{DoxyItemize}
\end{DoxyItemize}
\end{DoxyEnumerate}

{\bfseries General Information\-:}

List of potential energy terms compatible with Tinker-\/\-H\-P\-: Info about output of trajectories Info about free energy simulations, with specifics about alchemical ones through Lambda-\/\-A\-B\-F and related methodologies 